
%\begin{figure*}[t]
%	\centering
%	\includegraphics[width=0.9\linewidth, trim=18 18 18 18,clip]{fig/referring_segmentation_paradigms.pdf}
%	\caption{Comparison of referring segmentation paradigms. Existing systems expose segmentation through fixed token embeddings, polygon coordinates, or patch-level decisions. PLAS reads a grid of readout tokens in parallel, maps them to a spatially organized mask latent, aligns it with a teacher learned from ground-truth masks, and decodes it through a dense-prompt adapter and an image-conditioned mask decoder.}
%	\label{fig:paradigms}
%\end{figure*}
\section{Introduction}
\label{sec:intro}

Referring image segmentation (RIS) aims to identify
objects described by natural-language expressions
and delineate their precise pixel-level masks.
This task requires both semantic understanding
of image--text relationships and accurate
reconstruction of object geometry
\cite{Yu2016ModelingContext,Mao2016RefCOCOg,Liu2023GRES}.
Recently, multimodal large language models (MLLMs)
have demonstrated remarkable capabilities in
visual-language understanding and reasoning,
offering new opportunities for RIS
\cite{Liu2023LLaVA,Bai2025Qwen3VL}.
However, effectively transforming rich multimodal
representations within MLLMs into spatially
precise segmentation features remains challenging.

Existing MLLM-based segmentation methods commonly
connect multimodal representations to mask prediction
through different output interfaces, as summarized
in \cref{tab:paradigms}. Methods such as
LISA~\cite{Lai2024LISA} and
GSVA~\cite{Xia2024GSVA} employ segmentation-token
embeddings, while PixelLM~\cite{Ren2024PixelLM}
introduces a multi-scale segmentation codebook and
GLaMM~\cite{Rasheed2024GLaMM} combines segmentation
tokens with image-conditioned mask decoding.
Other methods, including
SAM4MLLM~\cite{Chen2024SAM4MLLM},
Text4Seg~\cite{Lan2025Text4Seg}, and
STAMP~\cite{Liu2026STAMP},
employ localization cues or patch-level predictions.
Although effective, these approaches rely on
different mechanisms to bridge multimodal understanding
and pixel-level segmentation.
Rather than focusing solely on output token counts
or decoding steps, we investigate how to extract
spatially organized segmentation features from
MLLM hidden states and explicitly align them with
fine-grained mask geometry before decoding.

In fact, MLLMs establish rich visual--language interactions
across multiple transformer layers, yet their hidden
states are not inherently optimized for pixel-level
segmentation. To address this challenge, we propose
\textbf{PLAS} (Progressive Latent Alignment with
Parallel Mask-Latent Prediction), an efficient
RIS framework built upon learnable multimodal
feature readout and progressive mask-latent
alignment. Specifically, we introduce a set of
learnable spatial readout tokens into the MLLM,
enabling them to interact with visual and textual
representations through the model's native
attention mechanism. By integrating multi-layer
hidden-state fusion, spatial refinement, and
visual-grid features, PLAS extracts spatially
organized, segmentation-aware representations
in parallel within a single MLLM forward pass.
The extracted features are subsequently mapped
into mask latents and transformed into dense
prompts for image-conditioned SAM decoding.

Nevertheless, parallel multimodal readout alone
does not guarantee the geometric precision
required for pixel-level segmentation.
We therefore develop a
\textbf{progressive teacher--student learning
	framework} to bridge the gap between multimodal
semantic understanding and mask geometry.
First, a mask-latent teacher learns
geometry-aware representations from ground-truth
masks through both image-free reconstruction
and image-conditioned SAM reconstruction.
Next, the student explicitly aligns its
spatially adapted parallel readout features
with the teacher-derived mask latents,
progressively acquiring mask-level geometric
knowledge.
Finally, image-conditioned task adaptation
transforms the predicted mask latents into
dense prompts for SAM-based decoding.
This progressive learning strategy establishes
a direct connection between multimodal feature
extraction and accurate pixel-level segmentation.

Our contributions are as follows:
\begin{itemize}
    \item We propose an \textbf{in-MLLM spatial readout
		architecture} with 256 learnable grid tokens,
	multi-layer hidden-state fusion, spatial refinement,
	and visual-grid injection, enabling one-pass extraction
	of segmentation-aware features.
	
	\item We construct a \textbf{geometry-aware mask-latent
		teacher} from ground-truth masks through image-free
	and image-conditioned reconstruction, providing a dense
	geometric target for multimodal feature learning.
	
	\item We develop a \textbf{progressive latent-to-mask
		decoding scheme}: the student is first warmed up by
	explicit latent alignment and then adapted with
	image-conditioned SAM; the predicted latent is converted
	into a dense prompt for final mask decoding.
	
	\item Experiments show a strong
	\textbf{accuracy--efficiency trade-off}: \textbf{85.7\%}
	cIoU on RefCOCO val at a model-side latency of \textbf{0.43 s/image}, with reported latency ratios corresponding to \textbf{1.63$\times$} and \textbf{2.49$\times$} faster inference than the fastest compared method and a method with comparable segmentation accuracy, respectively.
\end{itemize}
